Runtime guardrails for tool-using language-model agents must distinguish authorization, execution, and task completion. We evaluate authorization semantics and security--utility trade-offs after correcting an oracle monitor that compared complete action arguments against partial target annotations. Historical labels are retained without reinterpreting isolation rejections or failures as executed harmful effects.

A corrected DeepSeek Pro study comprises 900 authorization assignments over 180 fixed cases. Both runtime configurations use ordinary prompts. AgentDojo benign-tool completion rises from 2/30 under compiled authorization to 13/30 under exact-target policy and 11/30 under exact-target full modules. The corresponding M3-v2 differences are small and uncertain. Zero security events across 490 attack assignments do not establish effectiveness because this study has no unguarded control.

A separate 480-assignment contemporaneous calibration supplies that control. On AgentDojo, policy and full modules reduce benign completion from 17/30 to 2/30, a paired difference of \ensuremath{-}50 percentage points (95\% CI \ensuremath{-}66.83 to \ensuremath{-}26.80). Security results change from 3/30 to 0/30, but the paired interval includes zero. M3-v2 controls are already at zero. Sixteen synthetic Pro/Flash runs verify bounded private-data and missing-classification behavior, not general end-to-end protection. Missing outcomes and failures remain explicit.

The findings support measurement correction and disclosure of authorization burden, not full-module superiority. Oracle assistance, overlapping cohorts, adapted baselines, sparse security events, and incomplete independent-model validation limit generalization.
